
\documentclass{article}
\usepackage{colortbl}
\usepackage{makecell}
\usepackage{multirow}
\usepackage{supertabular}

\begin{document}

\newcounter{utterance}

\centering \large Interaction Transcript for game `cladder', experiment `full\_v1.5\_default', episode 4597 with score{-}Qwen3.5{-}9B{-}sp.
\vspace{24pt}

{ \footnotesize  \setcounter{utterance}{1}
\setlength{\tabcolsep}{0pt}
\begin{supertabular}{c@{$\;$}|p{.15\linewidth}@{}p{.15\linewidth}p{.15\linewidth}p{.15\linewidth}p{.15\linewidth}p{.15\linewidth}}
    \# & \multicolumn{2}{c}{Player} && \multicolumn{2}{c}{Game Master} \\
    \hline

    \theutterance \stepcounter{utterance}  
    & & & \multicolumn{4}{p{0.6\linewidth}}{
        \cellcolor[rgb]{0.9,0.9,0.9}{
            \makecell[{{p{\linewidth}}}]{
                \texttt{\tiny{[P1$\langle$GM]}}
                \texttt{You are an expert in causal inference. The following question is not a typical commonsense query, but rather a meticulously designed question created by a professor specializing in causal inference, intended to assess the students' mastery of the course content.} \\
\texttt{Imagine a self{-}contained, hypothetical world with only the following conditions, and without any unmentioned factors or causal relationships: Having a brother has a direct effect on wife and alarm clock. Wife has a direct effect on alarm clock. For people who do not have a brother and wives that don't set the alarm, the probability of ringing alarm is 88\%. For people who do not have a brother and wives that set the alarm, the probability of ringing alarm is 45\%. For people who have a brother and wives that don't set the alarm, the probability of ringing alarm is 58\%. For people who have a brother and wives that set the alarm, the probability of ringing alarm is 13\%. For people who do not have a brother, the probability of alarm set by wife is 9\%. For people who have a brother, the probability of alarm set by wife is 64\%. Does having a brother negatively affect alarm clock through wife?} \\
\texttt{Start your answer with 'Yes' or 'No', followed by additional reasoning or evidence to support your explanation.} \\
            }
        }
    }
     \\ \\

    \theutterance \stepcounter{utterance}  
    & \multicolumn{4}{p{0.6\linewidth}}{
        \cellcolor[rgb]{0.561,0.737,0.561}{
            \makecell[{{p{\linewidth}}}]{
                \texttt{\tiny{[P1$\rangle$GM]}}
                \texttt{No} \\
\\ 
\texttt{**Reasoning:**} \\
\\ 
\texttt{To determine if having a brother negatively affects the alarm clock *through* the wife, we must isolate the **indirect effect** of "Having a brother" on "Alarm ringing" via the mediator "Wife sets alarm."} \\
\\ 
\texttt{In this causal structure:} \\
\texttt{*   \$X\$: Having a brother (1 = Yes, 0 = No)} \\
\texttt{*   \$M\$: Wife sets alarm (1 = Yes, 0 = No)} \\
\texttt{*   \$Y\$: Alarm rings (1 = Yes, 0 = No)} \\
\\ 
\texttt{The indirect effect is calculated as the difference in the probability of the outcome (\$Y\$) when the mediator (\$M\$) is set to 1 versus 0, averaged over the distribution of the treatment (\$X\$), or more simply in this deterministic{-}path context, by comparing the conditional probabilities of \$Y\$ given \$M=1\$ and \$M=0\$ while holding \$X\$ constant.} \\
\\ 
\texttt{Let's look at the conditional probabilities of the alarm ringing provided in the prompt:} \\
\texttt{1.  **For people who do not have a brother (\$X=0\$):**} \\
\texttt{    *   If Wife sets alarm (\$M=1\$): \$P(Y=1 | X=0, M=1) = 45\textbackslash{}\%\$} \\
\texttt{    *   If Wife does not set alarm (\$M=0\$): \$P(Y=1 | X=0, M=0) =} \\
            }
        }
    }
    & & \\ \\

    \theutterance \stepcounter{utterance}  
    & & & \multicolumn{2}{p{0.3\linewidth}}{
        \cellcolor[rgb]{0.9,0.9,0.9}{
            \makecell[{{p{\linewidth}}}]{
                \texttt{\tiny{[GM$|$GM]}}
                \texttt{No} \\
            }
        }
    }
    & & \\ \\

    \theutterance \stepcounter{utterance}  
    & & & \multicolumn{2}{p{0.3\linewidth}}{
        \cellcolor[rgb]{0.9,0.9,0.9}{
            \makecell[{{p{\linewidth}}}]{
                \texttt{\tiny{[GM$|$GM]}}
                \texttt{yes} \\
            }
        }
    }
    & & \\ \\

    \theutterance \stepcounter{utterance}  
    & & & \multicolumn{2}{p{0.3\linewidth}}{
        \cellcolor[rgb]{0.9,0.9,0.9}{
            \makecell[{{p{\linewidth}}}]{
                \texttt{\tiny{[GM$|$GM]}}
                \texttt{game\_result = LOSE} \\
            }
        }
    }
    & & \\ \\

\end{supertabular}
}

\end{document}
